% docs/manual/user/chapters/03-config-structure.tex
% 第 3 章：配置文件结构

\chapter{配置文件结构}

第 2 章用 \cli{openbench init} 生成了一份最小可运行配置。本章详解 \openbench{} 的 YAML 配置文件格式：顶层四块如何组织、必填与可选字段如何配合默认值、以及两个高级语法（\yamlkey{_defaults} 与 \texttt{!include}）如何把大型多模型配置拆成可维护的多个文件。

完整字段参考见附录 A，本章重点解释结构与设计意图。

\section{顶层四块}

一个最小合法的 \file{openbench.yaml} 必须包含四个顶层 key：

\begin{minted}{yaml}
project:        # 项目元信息、时空范围、运行参数
evaluation:     # 要评估哪些变量
reference:      # 每个变量用哪个参考数据集
simulation:     # 要评估哪些模型（可多个）
\end{minted}

四个之外，还有 5 个常用可选块：

\begin{minted}{yaml}
metrics:        # list[str]，None 表示用全部 25+ 个 metric
scores:         # list[str]，None 表示用全部 normalized score
comparison:     # 多模型对比开关与项
statistics:     # 统计分析开关与项
options:        # （已废弃）旧的 v2.0 选项；自动迁移到 project
\end{minted}

\noteinline{加载器会把已废弃的 \yamlkey{options} 块自动迁移到 \yamlkey{project}，并打印 deprecation 警告。新配置请直接写到 \yamlkey{project} 下。}

\section{project 块}

\yamlkey{project} 是配置中字段最多的一块（25+ 字段），分为五组：

\subsection{身份标识（必填）}

\begin{minted}{yaml}
project:
  name: my-evaluation       # [必填] 项目名，用作输出目录子文件夹名
  output_dir: ./output      # [必填] 评估结果根目录（相对路径以工作目录为基）
  years: [2010, 2014]       # [必填] 评估时段 [起始年, 结束年]，含两端
\end{minted}

\warninline{\yamlkey{years} 必须是恰好 2 个整数的列表（不是 \texttt{[2010, 2011, 2012, 2013, 2014]}）。loader 会校验长度与类型并报错。}

\subsection{时空范围（可选，默认全球）}

\begin{minted}{yaml}
project:
  lat_range: [-90.0, 90.0]   # 默认全球
  lon_range: [-180.0, 180.0] # 默认全球
  min_year_threshold: 3      # 模型与参考重叠不足 N 年时跳过该变量
\end{minted}

\subsection{目标分辨率（可选，影响 evaluation/comparison/statistics 全部阶段）}

\begin{minted}{yaml}
project:
  tim_res: Month     # Month | Day | Hour | Year（不写则按解析规则推导）
  grid_res: 0.5      # 目标空间分辨率（度），不写则按解析规则推导
  timezone: 0        # 时区偏移（小时）
  weight: area       # 空间加权：none | area | mass（不写默认 area；详见第 7 章）
\end{minted}

\paragraph{tim\_res / grid\_res 不写时如何决定}

\modname{config.resolver.derive_target_resolution_context} 按下列优先级：

\begin{enumerate}
  \item \yamlkey{project.tim_res} / \yamlkey{project.grid_res} 已写 → 直接用
  \item 全部 simulation 的 \yamlkey{tim_res} / \yamlkey{grid_res} 一致 → 用这个共享值
  \item 各 simulation 的值不一致 → \texttt{ConfigError}：\emph{"Reference resolution is ambiguous across simulations: simulation grid\_res values differ; set project.grid\_res explicitly"}
\end{enumerate}

如果 1、2 都不命中（项目空、所有 sim 都没写），\modname{config.adapter} 在传给引擎前会再兜底：\yamlkey{tim_res} → \texttt{Month}，\yamlkey{grid_res} → \texttt{0.5}。但这是\emph{兜底而非默认}——实际工作流应该走优先级 1 或 2。

\paragraph{重网格的代价}

把所有 sim/ref 数据重采样到 \yamlkey{grid_res} 是 \openbench{} 在多模型/多分辨率混合评估时统一口径的关键。但 0.25° 全球 monthly + 30 个变量 + 6 个 sim 的重网格量大约 200--400 GB 中间态。如果几乎所有数据都已经在同一分辨率，写 \yamlkey{project.grid_res} 与原生匹配可避免无谓重采样。

\begin{tipBox}
\yamlkey{tim_res} 同样按这套规则推导。频率串支持 \texttt{Year} / \texttt{Month} / \texttt{Day} / \texttt{Hour} 以及 \texttt{3hr} / \texttt{6hr} 这种"数+单位"形式（\modname{data.processing.\_normalize\_frequency}）。
\end{tipBox}

\subsection{运行参数}

\begin{minted}{yaml}
project:
  num_cores: 8                  # 并行核数；省略则自动检测
  unified_mask: true            # 跨 sim 统一 NaN 掩码（保证比较公平）
  generate_report: true         # 生成 HTML/PDF 总结
  time_alignment: intersection  # intersection | per_pair | strict
\end{minted}

三种 \yamlkey{time_alignment} 策略：

\begin{itemize}
  \item \texttt{intersection}（默认）：所有 sim 与 ref 的共同时间窗口。安全保守
  \item \texttt{per_pair}：每对 sim×ref 各自的重叠区间。最大化数据利用，但跨对不可比
  \item \texttt{strict}：要求时间戳完全一致，否则报错
\end{itemize}

详见第 5 章。

\subsection{Groupby 分析}

\begin{minted}{yaml}
project:
  IGBP_groupby: false          # 按 IGBP 土地覆盖类型分组
  PFT_groupby: false           # 按 Plant Functional Type 分组
  climate_zone_groupby: false  # 按 Köppen 气候带分组
\end{minted}

\subsection{进阶开关}

\begin{minted}{yaml}
project:
  debug_mode: false        # 详细日志
  only_drawing: false      # 跳过计算，只重绘已有结果
  force: false             # 忽略增量缓存，全量重算
  strict_reference: false  # 把"低置信度元数据"警告升级为错误
\end{minted}

\section{evaluation 块}

最简单的一块，只有一个必填字段：

\begin{minted}{yaml}
evaluation:
  variables:
    - Evapotranspiration
    - Total_Runoff
    - Latent_Heat
\end{minted}

\yamlkey{variables} 必须是非空列表。变量名必须是 \openbench{} 已知的标准命名（驼峰式，下划线分隔），全部可用变量见第 4 章与附录 C。

\section{reference 块}

\yamlkey{reference} 把 \yamlkey{evaluation.variables} 中的每个变量映射到一个具体的 reference 数据集名：

\begin{minted}{yaml}
reference:
  data_root: /data/openbench/refs   # （可选）共享根目录
  Evapotranspiration: GLEAM_v4.2a_LowRes
  Total_Runoff: GRUN_ENSEMBLE_LowRes
  Latent_Heat: ERA5LAND_LowRes
\end{minted}

\warninline{\textbf{语法陷阱}：reference 必须是 \emph{flat} var→source 映射，不要写嵌套的 \texttt{sources:} 子块。早期 \cli{openbench init} 曾经写出嵌套形式导致 \cli{check} 失败 —— 现已修复，但旧示例文档若你看到，请按上面的 flat 形式写。}

\paragraph{每变量多个 reference}

同一个变量可以同时配多个 reference，\openbench{} 对每个 (sim × ref) 跑独立评估。三种语法等价：

\begin{minted}{yaml}
reference:
  # 1. List 形式（推荐，最清晰）
  Land_Surface_Temperature:
    - MODIST
    - TRIMS_LST

  # 2. 逗号分隔字符串（loader 会自动拆，方便从 v2.x namelist 直接迁移）
  Evapotranspiration: "GLEAM_v4.2a_LowRes, FLUXCOM-X-BASE_LowRes"

  # 3. 单 reference 仍是字符串（绝大多数变量这么写）
  Total_Runoff: GRUN_ENSEMBLE_LowRes
\end{minted}

详见第 4 章 §"每变量多 reference：多源对比"。

\subsection{data\_root：共享数据根}

\yamlkey{data_root} 是可选字段，不写时每个数据集用自己的 \yamlkey{root_dir}（registry 默认路径）。写了之后，grid 类数据集会用 \yamlkey{data_root} 替换 \yamlkey{root_dir}（station 类不受影响）。这在多用户共享的 HPC 环境很有用。

\subsection{自动分辨率后缀解析}

\yamlkey{reference} 里写 \texttt{GLEAM\_v4.2a} 时，\openbench{} 会根据 \yamlkey{project.grid_res} 自动加后缀：

\begin{itemize}
  \item \yamlkey{grid_res} >= 0.5 且 <= 1.0 → 优先尝试 \texttt{\_LowRes}
  \item \yamlkey{grid_res} = 0.25 → 优先尝试 \texttt{\_MidRes}
  \item 找不到对应分辨率版本时，回退到无后缀名
\end{itemize}

如果想绕过自动解析，直接写完整名（\texttt{GLEAM\_v4.2a\_MidRes}），就严格用那个分辨率。详见第 4 章。

\section{simulation 块}

\yamlkey{simulation} 是一个字典，key 是 \emph{标签}（你给这次评估自取的名字），value 是 simulation 元信息：

\begin{minted}{yaml}
simulation:
  CoLM2024_test:                  # 标签（任意，仅用于输出目录命名）
    model: CoLM2024               # [必填] 必须是已知 model profile 名
    root_dir: /data/colm2024_run  # [必填] 模型输出根目录

  CLM5_test:
    model: CLM5
    root_dir: /data/clm5_run
\end{minted}

每个 simulation 至少需要 \yamlkey{model} + \yamlkey{root_dir}。其余字段（\yamlkey{prefix} / \yamlkey{suffix} / \yamlkey{data_groupby} / \yamlkey{tim_res} / \yamlkey{grid_res} / \yamlkey{variables}）会从 model profile 继承默认值，可在此处局部覆盖。详见第 5 章。

\subsection{\yamlkey{_defaults}：simulation 块内的默认值合并}

如果你有多个 simulation 共享大量字段，可以用 \yamlkey{_defaults} 抽出公共部分：

\begin{minted}{yaml}
simulation:
  _defaults:
    data_type: grid
    grid_res: 0.5
    tim_res: Month
    prefix: case_

  Run_v1:
    model: CoLM2024
    root_dir: /data/v1
    suffix: _hist

  Run_v2:
    model: CoLM2024
    root_dir: /data/v2
    suffix: _hist_v2
\end{minted}

合并语义：

\begin{itemize}
  \item 每个 simulation 的有效字段 = \yamlkey{_defaults} 字段 + entry 自身字段
  \item entry 字段优先级高于 defaults
  \item \yamlkey{variables} 子字段做 \emph{深合并}（按变量名）：默认有 ET 和 LAI，entry 只覆盖 ET 时 LAI 仍在
  \item 其他字段做浅合并（覆盖式）
\end{itemize}

\warninline{\yamlkey{_defaults} 仅在 \yamlkey{simulation} 块内有效。在 \yamlkey{project} / \yamlkey{evaluation} / \yamlkey{reference} 等顶层块写 \yamlkey{_defaults} 不会被识别（loader 会把它当成异常 key 拒绝）。}

\section{!include：跨文件拆分}

当配置文件超过 200 行或多个团队共享部分配置时，可以把片段拆到独立 YAML 文件再 \texttt{!include} 进来。

\subsection{包含单个文件}

\begin{minted}{yaml}
# openbench.yaml
project: !include partials/project.yaml
evaluation:
  variables: !include partials/variables.yaml
reference: !include partials/refs.yaml
simulation: !include partials/sims.yaml
\end{minted}

每个被引入的文件就是一个普通 YAML：

\begin{minted}{yaml}
# partials/refs.yaml
data_root: /data/openbench/refs
Evapotranspiration: GLEAM_v4.2a_LowRes
Total_Runoff: GRUN_ENSEMBLE_LowRes
\end{minted}

路径相对于 \emph{包含 \texttt{!include} 的那个 YAML 文件所在目录}。

\subsection{Glob 模式包含多个文件}

\begin{minted}{yaml}
# openbench.yaml
simulation: !include simulations/*.yaml
\end{minted}

每个 \file{simulations/*.yaml} 应当是单条 simulation 字典：

\begin{minted}{yaml}
# simulations/colm2024.yaml
CoLM2024_test:
  model: CoLM2024
  root_dir: /data/colm2024_run
\end{minted}

匹配到的文件会按文件名排序后合并到同一个字典；同名 key 后加载者覆盖前加载者。

\subsection{递归 \texttt{!include}}

\openbench{} 的 loader 支持 \texttt{!include} 嵌套：被引入的文件里也可以有 \texttt{!include}，会递归展开。

\begin{tipBox}
\textbf{何时该用 !include}：当一个 YAML 区块在多个 project 间共享、或单个 project 的文件超过 ~150 行时；如果整体配置就 50 行，没必要拆。
\end{tipBox}

\section{完整最小示例}

把上面所有部分组合起来的最小示例：

\begin{minted}{yaml}
# openbench.yaml
project:
  name: demo
  output_dir: ./output
  years: [2010, 2014]
  tim_res: Month
  grid_res: 0.5
  num_cores: 8

evaluation:
  variables:
    - Evapotranspiration
    - Total_Runoff

reference:
  data_root: /data/openbench/refs
  Evapotranspiration: GLEAM_v4.2a
  Total_Runoff: GRUN_ENSEMBLE

simulation:
  CoLM2024:
    model: CoLM2024
    root_dir: /data/colm2024_run

comparison:
  enabled: true
\end{minted}

这份配置：
\begin{itemize}
  \item 评估 2010-2014 共 5 年的 ET 与 Runoff
  \item 用 GLEAM v4.2a + GRUN 作 reference，因 \yamlkey{grid_res: 0.5} 自动解析为 \texttt{\_LowRes} 后缀
  \item 用一个 CoLM2024 simulation
  \item 启用多模型对比（即便只有 1 个 sim 也会出对比图，仅展示自身）
\end{itemize}

\section{常见配置错误与诊断}

\subsection{Missing required field 错误}

\begin{minted}{text}
ConfigError: Missing required field: project.years
\end{minted}

确认 \yamlkey{project.name} / \yamlkey{output_dir} / \yamlkey{years} 三个必填字段都存在。

\subsection{years 列表长度不对}

\begin{minted}{text}
ConfigError: project.years must be a list of [start_year, end_year]
\end{minted}

\yamlkey{years} 是 \texttt{[2010, 2014]}，不是 \texttt{[2010, 2011, 2012, 2013, 2014]}。

\subsection{变量在 evaluation 列出但 reference 缺映射}

\begin{minted}{text}
ConfigError: Variable 'Latent_Heat' is in evaluation.variables but
has no entry in 'reference'.
Add: reference.Latent_Heat: <source_name>
\end{minted}

按提示在 \yamlkey{reference} 加一行映射即可。

\subsection{time\_alignment 拼错}

\begin{minted}{text}
ConfigError: project.time_alignment must be one of
{'intersection', 'per_pair', 'strict'}, got 'intersect'
\end{minted}

只有 3 个合法值，注意拼写与下划线。

\subsection{simulation 缺 _defaults 之外的条目}

\begin{minted}{text}
ConfigError: 'simulation' must contain at least one simulation entry
(not just _defaults)
\end{minted}

\yamlkey{simulation._defaults} 单独存在没意义，必须配至少一个具体 simulation。

\section{下一步}

下一章会详解：

\begin{itemize}
  \item 56 个评估变量按物理类别的全貌
  \item grid 与 station 数据形态在配置上的差异
  \item 67 个内置 reference 数据集如何浏览、下载、检查
  \item \cli{openbench data list} / \cli{download} / \cli{status} 用法
\end{itemize}
